\section*{Hoofdstuk \ref{ch:b1:curves} --- Vlakke krommen}

\begin{solution}{exo:b1:curves:1}
$M(-t) = (t^2, -t^3)$: spiegeling in de $x$-as; bestudeer $t \geq 0$.
Zowel $x' = 2t$ als $y' = 3t^2$ is $\geq 0$: de tak beweegt
naar rechts en omhoog, van $(0,0)$ naar oneindig. Bij $t = 0$ verdwijnt
de snelheid; de ontwikkelingen $x = t^2$, $y = t^3$ tonen dat de
raaklijn de $x$-as is ($y/x = t \to 0$) met $y$ dat van teken wisselt
terwijl $x \geq 0$: een \emph{\omterm{ex:b1:curves:astroid}{keerpunt}} dat naar links wijst. De kromme
is de semikubische parabool $y^2 = x^3$.
\end{solution}

\begin{solution}{exo:b1:curves:2}
$M(t + 2\pi) = M(t) + (2\pi, 0)$: de kromme herhaalt zich, verschoven met
één wielomtrek; bestudeer één periode. $x'(t) = 1 - \cos t \geq
0$, $y'(t) = \sin t$: de boog stijgt op $\intoo{0}{\pi}$, daalt op
$\intoo{\pi}{2\pi}$, met een top bij $(\pi, 2)$. Singuliere punten
waar $x' = y' = 0$: $t \in 2\pi\Z$, op de grond. Nabij $t = 0$:
\[
x(t) = \frac{t^3}{6} + o(t^3),
\qquad
y(t) = \frac{t^2}{2} + o(t^2):
\]
$x$ wisselt van teken, $y \geq 0$, en $\frac{x}{y^{3/2}}$ begrensd: de
raaklijn is \emph{verticaal} (richting $(0,1)$), een \omterm{ex:b1:curves:astroid}{keerpunt} waar
het gevolgde punt momentaan nulsnelheid heeft --- de fysische signatuur
van rollen zonder slippen.
\end{solution}

\begin{solution}{exo:b1:curves:3}
Bij $t = \frac\pi4$: $M = \bigl(\frac{\sqrt2}{4},
\frac{\sqrt2}{4}\bigr)$, raakrichting $(-\cos t, \sin t) =
\frac{1}{\sqrt2}(-1, 1)$ (\cref{ex:b1:curves:astroid}). Raaklijn:
$y - \frac{\sqrt2}{4} = -(x - \frac{\sqrt2}{4})$, d.w.z.\ $x +
y = \frac{\sqrt2}{2}$. Snijpunten met de assen: $\bigl(\frac{\sqrt2}{2},
0\bigr)$ en $\bigl(0, \frac{\sqrt2}{2}\bigr)$; het segment ertussen
heeft \omterm{def:b1:curves:arclength}{lengte} $\sqrt{\frac12 + \frac12} = 1$.

(Algemeen punt $t$: de raaklijn in $(\cos^3 t, \sin^3 t)$ snijdt de
assen in $(\cos t, 0)$ en $(0, \sin t)$ --- controleer dat de rechte
door deze punten richting $(-\cos t, \sin t)$ heeft en door
$M(t)$ gaat --- en het afgesneden segment heeft \omterm{def:b1:curves:arclength}{lengte} $\sqrt{\cos^2 t +
\sin^2 t} = 1$.)
\end{solution}

\begin{solution}{exo:b1:curves:4}
$r = \cos 2\theta$: periode $\pi$ in $\theta$, symmetrisch in beide
assen; $r$ verdwijnt bij $\theta = \pm\frac\pi4$ (raaklijnen in de
oorsprong langs de diagonalen) en is negatief voor $\theta \in
\intoo{\frac\pi4}{\frac{3\pi}{4}}$, waar de punten op de
tegenoverliggende straal worden afgezet --- wat vier blaadjes langs de
asrichtingen $\theta = 0, \frac\pi2, \pi, \frac{3\pi}{2}$ oplevert, elk
met maximale straal $1$.

$r\cos\theta = 1$ is de vergelijking $x = 1$: de \omterm{def:b1:curves:polar}{poolkromme} $r =
\frac{1}{\cos\theta}$ is de verticale rechte $x = 1$ (één keer doorlopen
voor $\theta \in \intoo{-\frac\pi2}{\frac\pi2}$).
\end{solution}

\begin{solution}{exo:b1:curves:5}
$r(\theta) = 1 + 2\cos\theta$ verdwijnt voor $\cos\theta = -\frac12$:
$\theta = \pm\frac{2\pi}{3}$. Symmetrie in de $x$-as; bestudeer
$\theta \in \intcc{0}{\pi}$. Voor $\theta \in
\intco{0}{\frac{2\pi}{3}}$: $r > 0$, dalend van $3$ naar $0$:
buitenboog, de oorsprong binnengaand rakend aan de straal $\theta =
\frac{2\pi}{3}$. Voor $\theta \in \intoc{\frac{2\pi}{3}}{\pi}$: $r <
0$: de punten $r\vec u(\theta)$ liggen op de tegenoverliggende straal
(hoek $\theta - \pi \in \intoc{-\frac\pi3}{0}$), met afstand $\abs r$
groeiend van $0$ naar $1$: dit tekent een kleine \emph{binnenlus}
door de oorsprong. Bij $\theta = \pi$, $r = -1$ en het punt is
$-\vec u(\pi) = (1, 0)$: de lus sluit op de positieve $x$-as.
Schets: een grote hartachtige buitenkromme
met maximaal bereik $3$ bij $\theta = 0$, plus een lus erbinnen
door $O$ en $(1,0)$.
\end{solution}

\begin{solution}{exo:b1:curves:6}
Vervolledig de kwadraten:
\[
9(x^2 - 4x) + 25(y^2 - 2y) = 164
\iff 9(x-2)^2 + 25(y-1)^2 = 164 + 36 + 25 = 225 .
\]
Delen door $225$: $\frac{(x-2)^2}{25} + \frac{(y-1)^2}{9} = 1$: een
ellips met middelpunt $(2, 1)$, halve assen $a = 5$ (horizontaal), $b =
3$; $c = \sqrt{25 - 9} = 4$: brandpunten $(2 \pm 4,\, 1) = (-2, 1)$ en
$(6, 1)$; excentriciteit $e = \frac45$.
\end{solution}

\begin{solution}{exo:b1:curves:7}
Brandpunt in de oorsprong, richtlijn $D: x = d$ met $d = \frac pe > 0$.
Voor een punt $M = (r\cos\theta, r\sin\theta)$ met $r > 0$:
\[
MF = r,
\qquad
d(M, D) = \abs{d - r\cos\theta},
\]
en de kegelsnedevoorwaarde $MF = e\,d(M, D)$, in het regime waar $M$
aan de brandpuntzijde van de richtlijn ligt ($r\cos\theta < d$), luidt $r
= e(d - r\cos\theta)$, d.w.z.
\[
r(1 + e\cos\theta) = ed = p,
\qquad
r = \frac{p}{1 + e\cos\theta} .
\]
Voor $e < 1$ verdwijnt de noemer nooit: alle $\theta$ toegelaten
(ellips). Voor $e = 1$: $\theta \neq \pi$ (parabool, \omterm{def:b1:topology:open}{open} naar
de verre zijde van de richtlijn). Voor $e > 1$: nodig is $1 + e\cos\theta > 0$,
d.w.z.\ $\theta \in \intoo{-\theta_0}{\theta_0}$ met $\theta_0 =
\arccos\bigl(-\frac1e\bigr)$: één \emph{tak} van de hyperbool
(de andere tak stemt overeen met de tekenkeuze $r < 0$, of met het
tweede brandpunt).
\end{solution}

\begin{solution}{exo:b1:curves:8}
\emph{Tuinman:} bevestig een touw van \omterm{def:b1:curves:arclength}{lengte} $2a$ aan twee palen $F,
F'$ en houd het strak met het tekenende punt $M$: de beperking is
precies $MF + MF' = 2a$, dus de getekende kromme is de ellips.

\emph{Spiegeleigenschap:} zij $t \mapsto M(t)$ een reguliere
parametrisatie en $u(t) = \frac{M(t) - F}{\norm{M(t) - F}}$,
$u'(t)$ de analoge eenheidsvector naar $F'$. Door
$\norm{M - F} + \norm{M - F'} = 2a$ te differentiëren, met behulp van
$\frac{\dd}{\dd t}\norm{M - F} = \bigl\langle \frac{M - F}{\norm{M-F}},\,
M'\bigr\rangle$ (kettingregel op $\sqrt{\langle\cdot,\cdot\rangle}$):
\[
\bigl\langle u + u',\, M'\bigr\rangle = 0 .
\]
Dus de raakrichting $M'$ is orthogonaal aan de bissectricerichting
$u + u'$ van de twee brandpuntstralen: de raaklijn maakt gelijke
hoeken met $MF$ en $MF'$. (Een lichtstraal vanuit één brandpunt
weerkaatst tegen de ellips naar het andere brandpunt.)
\end{solution}

\begin{solution}{exo:b1:curves:9}
\begin{enumerate}
  \item Domein: $\cos 2\theta \geq 0$: $\theta \in
        \intcc{-\frac\pi4}{\frac\pi4} \cup
        \intcc{\frac{3\pi}{4}}{\frac{5\pi}{4}}$. Symmetrieën: $\theta
        \mapsto -\theta$ ($x$-as) en $\theta \mapsto \pi -
        \theta$ ($y$-as): bestudeer $\intcc{0}{\frac\pi4}$; $r$
        daalt van $1$ naar $0$. Bij $\theta = \pm\frac\pi4$: $r =
        0$, raaklijnen in de oorsprong langs de diagonalen. De kromme is
        het $\infty$-symbool: twee symmetrische lussen die elkaar
        ontmoeten in $O$ en $(\pm 1, 0)$ bereiken.
  \item Met $F, F' = (\pm c, 0)$, $c = \frac{1}{\sqrt 2}$, en $M
        = (r\cos\theta, r\sin\theta)$:
        \[
        MF^2\, MF'^2
        = \bigl((r\cos\theta - c)^2 + r^2\sin^2\theta\bigr)
        \bigl((r\cos\theta + c)^2 + r^2\sin^2\theta\bigr)
        = (r^2 + c^2)^2 - (2rc\cos\theta)^2 ,
        \]
        volgens de identiteit $(A - B)(A + B) = A^2 - B^2$ met $A = r^2
        + c^2$, $B = 2rc\cos\theta$. Met $c^2 = \frac12$:
        \[
        MF^2 MF'^2 = \Bigl(r^2 + \frac12\Bigr)^2 - 2r^2\cos^2\theta
        = r^4 + r^2\bigl(1 - 2\cos^2\theta\bigr) + \frac14
        = r^4 - r^2\cos 2\theta + \frac14 .
        \]
        Op de lemniscaat, $r^2 = \cos 2\theta$: de eerste twee termen
        vallen weg, wat $MF^2MF'^2 = \frac14$ overlaat, d.w.z.\ $MF \cdot MF' =
        \frac12$. Omgekeerd toont de berekening achterwaarts gelezen dat
        de meetkundige-plaatsvergelijking $MF\,MF' = \frac12$ gelijk is aan $r^2 = \cos
        2\theta$ (voor $r \neq 0$; en $O$ voldoet aan beide).
\end{enumerate}
\end{solution}

\begin{solution}{exo:b1:curves:10}
$r = 1 + \cos\theta$, $r' = -\sin\theta$:
\[
r^2 + r'^2 = 1 + 2\cos\theta + \cos^2\theta + \sin^2\theta
= 2 + 2\cos\theta = 4\cos^2\frac\theta2 ,
\]
dus $\sqrt{r^2 + r'^2} = 2\abs{\cos\frac\theta2}$ en, door de
symmetrie in de $x$-as,
\[
L = \int_0^{2\pi} 2\Bigl|\cos\frac\theta2\Bigr|\,\dd\theta
= 2\int_0^{\pi} 2\cos\frac\theta2\,\dd\theta
= 2\Bigl[4\sin\frac\theta2\Bigr]_0^{\pi} = 8 .
\]
Nog een algebraïsche, $\pi$-vrije omtrek.
\end{solution}

\begin{solution}{exo:b1:curves:11}
Het punt $(a\cos t, b\sin t)$ voldoet aan de vergelijking:
$\cos^2 t + \sin^2 t = 1$. De normaalvector van de rechte is
$\bigl(\frac{\cos t}a, \frac{\sin t}b\bigr)$, en zijn product
met de snelheid $(-a\sin t, b\cos t)$ is $-\sin t\cos t +
\sin t\cos t = 0$: de rechte gaat door het punt met de
raakrichting --- ze is de raaklijn. Voor $(x_0, y_0)$ op de
ellips, schrijf $\cos t = \frac{x_0}a$, $\sin t = \frac{y_0}b$ en
substitueer:
\[
\frac{x\,x_0}{a^2} + \frac{y\,y_0}{b^2} = 1 ,
\]
verkregen uit de ellipsvergelijking door $x^2 \mapsto
x\,x_0$ en $y^2 \mapsto y\,y_0$ te ``splitsen''.
\end{solution}

\begin{solution}{exo:b1:curves:12}
\begin{enumerate}
  \item $r' = k\eu^{k\theta}$, dus $\tan V = \frac{r}{r'} =
        \frac1k$: constant. De spiraal snijdt elke straal vanuit de
        oorsprong onder dezelfde hoek $V = \arctan\frac1k$.
  \item In complexe notatie is de spiraal $\{\eu^{k\theta}
        \eu^{\iu\theta} : \theta \in \R\}$. Roteren over $c$
        vermenigvuldigt met $\eu^{\iu c}$:
        \[
        \eu^{k\theta}\eu^{\iu(\theta + c)}
        = \eu^{-kc}\,\eu^{k(\theta + c)}\eu^{\iu(\theta+c)} ,
        \]
        en terwijl $\theta + c$ over $\R$ loopt, beschrijft dit
        $\eu^{-kc}$ maal de spiraal: rotatie $=$ schaling. Geen
        andere gladde kromme dan rechten en cirkels heeft deze
        eigenschap.
  \item $\sqrt{r^2 + r'^2} = \sqrt{1 + k^2}\,\eu^{k\theta}$, dus
        op $\intcc{A}{\theta_0}$ is de \omterm{def:b1:curves:arclength}{lengte}
        $\frac{\sqrt{1+k^2}}{k}\bigl(\eu^{k\theta_0} -
        \eu^{kA}\bigr)$, en als $A \to -\infty$:
        \[
        L = \frac{\sqrt{1 + k^2}}{k}\,\eu^{k\theta_0} < \infty :
        \]
        oneindig veel omwentelingen rond de oorsprong, eindige totale
        \omterm{def:b1:curves:arclength}{lengte} (de omwentelingen krimpen meetkundig).
\end{enumerate}
\end{solution}

\begin{solution}{pb:b1:curves:1}
\textbf{1.} Rollen zonder slippen betekent dat het contactpunt een
afstand heeft afgelegd gelijk aan de boog van het afgerolde wiel: na
draaien over $t$ ligt het middelpunt in $(Rt, R)$. Het gemarkeerde punt
ligt op de \omterm{def:b1:topology:closure}{rand} onder hoek $t$ \emph{achter} de neerwaartse verticale
(het wiel draait met de klok mee terwijl het vooruitgaat):
\[
M(t) = \Omega(t) + R(-\sin t, -\cos t)
= \bigl(R(t - \sin t),\ R(1 - \cos t)\bigr),
\]
wat correct is bij $t = 0$ ($M = (0,0)$) en bij $t = \pi$ ($M =
(\pi R, 2R)$, de top).

\textbf{2.} $f'(t) = R(1 - \cos t,\ \sin t)$ en
\[
\norm{f'}^2 = R^2\bigl((1-\cos t)^2 + \sin^2 t\bigr)
= 2R^2(1 - \cos t) = 4R^2\sin^2\frac t2 ,
\]
dus $\norm{f'} = 2R\abs{\sin\frac t2}$. Singuliere punten bij $t \in
2\pi\Z$: de \omterm{ex:b1:curves:astroid}{keerpunten} op de grond (\cref{exo:b1:curves:2}); de
top van de boog is $t = \pi$, waar de snelheid $2R$ maximaal is.

\textbf{3.} Halvehoekfactorisaties:
\[
f'(t) = 2R\sin\frac t2\,\Bigl(\sin\frac t2,\ \cos\frac t2\Bigr),
\quad
M - C = 2R\sin\frac t2\,\Bigl(-\cos\frac t2,\ \sin\frac t2\Bigr),
\]
\[
T - M = \bigl(R\sin t,\ R(1 + \cos t)\bigr)
= 2R\cos\frac t2\,\Bigl(\sin\frac t2,\ \cos\frac t2\Bigr).
\]
Voor $0 < t < 2\pi$, $\sin\frac t2 \neq 0$: $M - C$ is orthogonaal
aan de raakrichting $\bigl(\sin\frac t2, \cos\frac
t2\bigr)$ (hun product is $-\sin\frac t2\cos\frac t2 +
\sin\frac t2\cos\frac t2 = 0$), en $T - M$ is er evenwijdig mee.
De koorde naar de top van het wiel \emph{is} de raaklijn; de
koorde naar het contactpunt is de normaal.

\textbf{4.} Op elk ogenblik draait het wiel om zijn contactpunt (dat
punt heeft nulsnelheid: rollen zonder slippen),
dus elk star punt van het wiel beweegt orthogonaal op de rechte
die het met $C$ verbindt, met snelheid (hoeksnelheid $1$) gelijk aan
die afstand. Controle: $\norm{M - C} = 2R\abs{\sin\frac t2} =
\norm{f'(t)}$.

\textbf{5.} $2R\,y(t) = 2R\cdot 2R\sin^2\frac t2 =
4R^2\sin^2\frac t2 = \norm{f'(t)}^2$. De snelheid op hoogte $y$ is
$\sqrt{2R\,y}$ --- formeel de wet $v = \sqrt{2gh}$ van de vrije
val, met de grond die het plafond speelt; Deel IV maakt van deze
observatie uurwerk.

\textbf{6.} Volgens \cref{def:b1:curves:arclength} en vraag 2
($\sin\frac t2 \geq 0$ op $\intcc{0}{2\pi}$):
\[
L = \int_0^{2\pi} 2R\sin\frac t2\,\dd t
= 2R\Bigl[-2\cos\frac t2\Bigr]_0^{2\pi} = 2R(2 + 2) = 8R .
\]
Wrens stelling: precies vier diameters.

\textbf{7.} $s(t) = \int_0^t 2R\sin\frac u2\,\dd u =
4R\bigl(1 - \cos\frac t2\bigr)$; $s(2\pi) = 4R(1+1) = 8R$,
consistent.

\textbf{8.} $\sigma(t) = \abs{s(t) - s(\pi)} = \abs{4R(1 -
\cos\frac t2) - 4R} = 4R\abs{\cos\frac t2}$, dus $\sigma^2 =
16R^2\cos^2\frac t2$; en $2R - y = R(1 + \cos t) =
2R\cos^2\frac t2$, waaruit $8R(2R - y) = 16R^2\cos^2\frac t2 =
\sigma^2$.

\textbf{9.} Bij het \omterm{ex:b1:curves:astroid}{keerpunt} $t = 0$ (of $2\pi$): $\sigma = 4R$, de
helft van $8R$: de top halveert de boog. In beide berekeningen is de
snelheid $\abs{\text{trigonometrische veelterm in } t/2}$, waarvan de
primitieve opnieuw trigonometrisch is: de \omterm{def:b1:curves:arclength}{lengte} is een
verschil van \emph{waarden} van cosinussen --- rationale getallen
maal $R$ --- zonder cirkelboog te meten, dus geen $\pi$.
De cirkel zelf heeft constante snelheid, dus zijn \omterm{thm:b1:integration:def}{lengte-integraal}
levert de volledige intervallengte $2\pi$; de snelheid van de cycloïde
verdwijnt aan de uiteinden en integreert algebraïsch.

\textbf{10.} De boog wordt doorlopen met $x$ stijgend van $0$ tot
$2\pi R$ ($x' = R(1 - \cos t) \geq 0$, enkel verdwijnend in
geïsoleerde punten). Het substitueren van $x = x(t)$ in de
\omterm{thm:b1:integration:def}{oppervlakte-integraal} $\int_0^{2\pi R} y\,\dd x$ (\cref{thm:b1:integration:parts})
geeft $A = \int_0^{2\pi} y(t)\,x'(t)\,\dd t$.

\textbf{11.}
\[
A = \int_0^{2\pi} R(1 - \cos t)\cdot R(1 - \cos t)\,\dd t
= R^2\int_0^{2\pi}\bigl(1 - 2\cos t + \cos^2 t\bigr)\dd t
= R^2\Bigl(2\pi - 0 + \pi\Bigr) = 3\pi R^2 ,
\]
met behulp van $\int_0^{2\pi}\cos^2 = \pi$. Precies drie wieloppervlakten:
Galilei's balans zei ``ongeveer $3$''; Robervals berekening
zegt ``precies''.

\textbf{12.} Met $x = \cos^3 t$, $y = \sin^3 t$: $y\,x' =
-3\sin^4 t\cos^2 t$. Lineariseer:
\[
\sin^4 t\cos^2 t = \frac{1 - \cos 2t}{2}\cdot\frac{\sin^2
2t}{4} = \frac{\sin^2 2t}{8} - \frac{\sin^2 2t\cos 2t}{8} ,
\]
en over $\intcc{0}{2\pi}$: $\int\sin^2 2t = \pi$, $\int \sin^2
2t\cos 2t = \bigl[\frac{\sin^3 2t}{6}\bigr] = 0$. Dus $\oint
y\,\dd x = -3\cdot\frac\pi8$, en de ingesloten oppervlakte is
$\frac{3\pi}8$ (het teken registreert de tegenwijzerzin-
oriëntatie).

\textbf{13.} Voor $(\cos t, \sin t)$ met $t$ van $\pi$ naar $0$,
stijgt $x$ van $-1$ naar $1$ en
\[
\int_\pi^0 \sin t\cdot(-\sin t)\,\dd t
= \int_0^\pi \sin^2 t\,\dd t = \frac\pi2 :
\]
de formule geeft de oppervlakte van de bovenste halve schijf terug, zoals het moet.

\textbf{14.} $x' = R(1 + \cos t) = 2R\cos^2\frac t2$, $y' =
R\sin t = 2R\sin\frac t2\cos\frac t2$, dus $\norm{f'} =
2R\cos\frac t2$ (niet-negatief op $\intcc{-\pi}{\pi}$). Boog vanaf
de bodem: $s(t) = \int_0^t 2R\cos\frac u2\,\dd u =
4R\sin\frac t2$. Dan
\[
y = R(1 - \cos t) = 2R\sin^2\frac t2
= 2R\Bigl(\frac{s}{4R}\Bigr)^{\!2} = \frac{s^2}{8R} .
\]

\textbf{15.} Behoud van energie met $v = \frac{\dd s}{\dd\tau}$
en $y = \frac{s^2}{8R}$, $y_0 = \frac{s_0^2}{8R}$:
\[
\Bigl(\frac{\dd s}{\dd\tau}\Bigr)^{\!2} = 2g(y_0 - y)
= \frac{2g}{8R}\bigl(s_0^2 - s^2\bigr)
= \frac{g}{4R}\bigl(s_0^2 - s^2\bigr).
\]

\textbf{16.} Differentiëren naar $\tau$: $2s's'' =
-\frac{g}{4R}\,2ss'$, dus overal waar $s' \neq 0$ (dus overal
door \omterm{def:b1:continuity:continuous}{continuïteit}), $s'' = -\frac{g}{4R}s$: de harmonische oscillator.
Volgens \cref{thm:b1:diffeq:homogeneous2}, $s(\tau) =
A\cos\omega\tau + B\sin\omega\tau$ met $\omega =
\sqrt{g/(4R)}$; de beginvoorwaarden $s(0) = s_0$, $s'(0) = 0$
geven $s(\tau) = s_0\cos\omega\tau$.

\textbf{17.} De kraal bereikt de bodem wanneer $s = 0$, d.w.z.\ bij
$\omega\tau = \frac\pi2$:
\[
T_{\downarrow} = \frac{\pi}{2\omega}
= \frac\pi2\sqrt{\frac{4R}{g}} = \pi\sqrt{\frac Rg}\,,
\]
onafhankelijk van $s_0$: losgelaten van waar dan ook in de schaal, komen
alle kralen op hetzelfde ogenblik aan --- de tautochroon.

\textbf{18.} Door $s = s_0\cos\omega\tau$ te substitueren: het linkerlid
is $s_0^2\omega^2\sin^2\omega\tau$ en het rechterlid is
$\frac{g}{4R}s_0^2(1 - \cos^2\omega\tau) =
s_0^2\omega^2\sin^2\omega\tau$: exacte gelijkheid, dus er werd geen
onechte oplossing ingevoerd. Voor een cirkelboog is $y$ niet
evenredig met $s^2$ ($y = \ell(1 - \cos\frac s\ell) =
\frac{s^2}{2\ell} - \frac{s^4}{24\ell^3} + \dots$): de
herstelterm is slechts \emph{bij benadering} \omterm{def:b1:linmaps:def}{lineair}, zodat de
periode van een cirkelvormige slinger met de amplitude verschuift, terwijl de
periode van de cycloïde streng constant is.

\textbf{19.} $\pi\sqrt{R/g} = 1$ geeft $R = \frac{g}{\pi^2}
\approx \frac{9{,}81}{9{,}87} \approx 0{,}994$ m --- bijna precies
één meter. De volledige oscillatie duurt $\frac{2\pi}\omega =
2\pi\sqrt{4R/g}$, wat precies de kleinehoekformule
$2\pi\sqrt{\ell/g}$ is voor een slinger van \omterm{def:b1:curves:arclength}{lengte} $\ell = 4R$:
Huygens hing zijn slinger op aan cycloïdale wangen van precies
die verhouding.

\textbf{20.} Bij $t = \frac\pi2$, $R = 1$: $M = \bigl(\frac\pi2 -
1,\ 1\bigr)$, $T = \bigl(\frac\pi2,\ 2\bigr)$, dus $T - M = (1,
1)$. En $f'(\frac\pi2) = (1 - 0,\ 1) = (1, 1)$: de koorde naar
de top is de snelheid, zoals beloofd.

\textbf{21.} Voor de trochoïde, $x'(t) = R - d\cos t$ en $y'(t)
= d\sin t$. Als $d < R$: $x' \geq R - d > 0$, het punt gaat steeds
vooruit; de snelheid verdwijnt nooit (haar eerste component is
positief): geen singulier punt, een reguliere golf. Als $d > R$: $x'(0)
= R - d < 0 < R + d = x'(\pi)$, dus het punt beweegt achteruit
nabij de contactogenblikken en vooruit elders: de kromme
kruist zichzelf in lussen. Een punt op de flens van een spoorweg
wiel (onder de railkop, $d > R$) reist bij elke omwenteling achteruit.

\textbf{22.} De koorde van het \omterm{ex:b1:curves:astroid}{keerpunt} $(\pi R, 2R)$ naar de oorsprong
heeft \omterm{def:b1:curves:arclength}{lengte} $L = R\sqrt{\pi^2 + 4}$ en hellingshoek $\alpha$ met
$\sin\alpha = \frac{2R}{L}$. Glijden vanuit rust met constante
versnelling $g\sin\alpha$: $L = \frac12 g\sin\alpha\,T^2$, dus
\[
T = \sqrt{\frac{2L}{g\sin\alpha}} = \sqrt{\frac{2L^2}{2Rg}}
= \frac{L}{\sqrt{Rg}} = \sqrt{\pi^2 + 4}\,\sqrt{\frac Rg}
\approx 3{,}72\sqrt{\frac Rg}\,,
\]
tegenover $\pi\sqrt{R/g} \approx 3{,}14\sqrt{R/g}$ voor de cycloïde:
het gekromde pad is sneller. Het is in feite het snelst mogelijke
--- de brachistochroon --- een stelling van de variatie
rekening, buiten dit volume.

\textbf{23.} $\dfrac{\dd y}{\dd x} = \dfrac{y'}{x'} =
\dfrac{\sin t}{1 - \cos t} = \cot\dfrac t2$ (halvehoek
formules). Dan
\[
\frac{\dd^2y}{\dd x^2}
= \frac{\dd}{\dd t}\Bigl(\cot\frac t2\Bigr)\cdot\frac{1}{x'(t)}
= -\frac{1}{2\sin^2\frac t2}\cdot\frac{1}{2R\sin^2\frac t2}
= -\frac{1}{4R\sin^4\frac t2} < 0 :
\]
concaaf over de hele boog; als $t \to 0^+$ of $2\pi^-$ gaat de
helling $\cot\frac t2 \to \pm\infty$: verticale raaklijnen in de
\omterm{ex:b1:curves:astroid}{keerpunten}.

\textbf{24.} De koorderichting is $T - M \parallel
\bigl(\sin\frac t2, \cos\frac t2\bigr)$, die naar $(0, 1)$ streeft
als $t \to 0^{+}$: de raaklijn in het \omterm{ex:b1:curves:astroid}{keerpunt} is verticaal ---
herwonnen uit pure meetkunde, zonder Taylor-ontwikkeling.

\textbf{25.} (i) Wrens stelling, $L = 8R$; Robervals oppervlakte, $A =
3\pi R^2$ (Galilei's vermoede verhouding $3$); Huygens'
tautochroon, $T_\downarrow = \pi\sqrt{R/g}$. (ii) Elke
berekening liep op de halvehoekfactorisaties $f' =
2R\sin\frac t2(\sin\frac t2, \cos\frac t2)$ en hun schaal
analogon --- één identiteit die raaklijn, \omterm{def:b1:curves:arclength}{lengte} en klok
tegelijk aandreef. (iii) De intrinsieke relatie $y = s^2/8R$ zette de
energievergelijking om in $s'' = -\frac{g}{4R}s$, een lineaire vergelijking
met constante coëfficiënten waarvan de oplossingen exact
isochroon zijn. (iv) Deel I gebruikte de differentiaalrekening van
\cref{ch:b1:derivative}, Deel II--III de \omterm{thm:b1:integration:def}{integraal} van
\cref{ch:b1:integration}, Deel IV de differentiaalvergelijkingen van
\cref{ch:b1:diffeq} --- de cycloïde is het curriculum van dit boek
opgerold tot één kromme.
\end{solution}
